\uuid{lENH}
\titre{Résolution d'équation différentielle du premier ordre}
\auteur{bodin ; croizat ; sacre}
\organisation{exo7}
\licence{CC BY-NC-SA 4.0}
\datecreate{2022-03-08}
\niveau{L1}
\module{Analyse}
\chapitre{Equation différentielle}
\sousChapitre{Résolution d'équation différentielle du premier ordre}
\difficulte{3}
\theme{}
\exo7qcmid{hilisit.equations-differentielles.051}

\contenu{
\begin{question}
\qclassification{225.01}
Le graphique ci-dessous représente la courbe représentative d'une fonction \(f\) ainsi que sa tangente en un point \(A\). Cette fonction \(f\) est solution d'une des équations différentielles suivantes ; laquelle ?

\qimage{images/png/lENH-1.png}

\begin{answers}
	\bad{\(y'=2x\)}
	\bad{\(y'=y+1\)}
	\good{\(y'=2y+2\)}
	\bad{\(y'=2y-2\)}
\end{answers}
\begin{explanations}
On a \(f(0)=0\) et \(f'(0)=2\) puisqu'il s'agit de la pente de la tangente à la courbe de $f$ au point d'abscisse $x=0$. Ceci élimine toutes les réponses proposées sauf \(y'=2y+2\). De fait, la courbe représentative donnée est celle de \(x\mapsto  e^{2x}-1\) qui en est bien une solution.
\end{explanations}
\end{question}
}
